\scanpage{134}
\begin{center}- \mspage{115} -\end{center}\margin{IX}

enfin donnée sur $\eta$ un pro-$\ell$-groupe de Barsotti-Tate $X_\eta$ qui soit \uline{extension panachée} de $E_\eta$ par $F_\eta$, par quoi nous entendons qu'il est muni d'une filtration en trois crans (en tant que pro-$\ell$-groupe de Barsotti-Tate), donnant lieu à des extensions partielles munies d'isomorphismes avec $E_\eta$ et $F_\eta$ respectivement. Il s'ensuit que pour tout $n$, $X_\eta(n)$ est muni d'une structure d'extension panachée de $E(n)_\eta$ par $F(n)_\eta$. Nous pouvons lui faire correspondre alors un accouplement du type (9.4.1) :
\[
(9.5.4)\qquad c(X(n)_\eta):Q(n)\otimes Q^*(n)\longrightarrow\Lambda_n=\mathbb Z/\ell^{n+1}\mathbb Z,
\]
où les $Q^*(n)$ sont les groupes finis étales sur $S$ définis par
\[
P(n)\simeq D(Q^*(n)),\quad\text{i.e.}\quad Q^*(n)=D(P(n)),
\]
$D$ désignant encore la dualité de Cartier. Les $Q^*(n)$ forment un pro-$\ell$-groupe de façon évidente, soit $Q^*$. Nous laissons au lecteur le soin de vérifier que pour $n$ variable, les homomorphismes (9.5.4) définissent un homomorphisme de systèmes projectifs, i.e. un homomorphisme de faisceaux $\ell$-adiques constants tordus sur $S$ :
\[
(9.5.5)\qquad c(X_\eta):Q\otimes Q^*\longrightarrow\mathbb Z_{\ell S}.
\]

Signalons en passant le résultat suivant (qui ne sera pas utilisé dans la suite) :

\uline{Proposition} 9.5.6. \uline{Les données} $P,Q,R,E,F$ \uline{sont les mêmes que ci-dessus}, \uline{sauf qu'il est inutile de supposer le trait} $S$ \uline{hensélien}. \uline{On considère le foncteur}
\[
(9.5.6.1)\qquad X\longmapsto X_\eta
\]
\uline{de la catégorie des extensions panachées de Barsotti-Tate de} $E$ \uline{par} $F$

\margin{\unsure{433|432}}\margin{433}\margin{133}
